\section{Problem and Design Opportunity}
\label{sec:problem}

The design opportunity arises from separating three decisions that request
routing often couples: \emph{placement}, \emph{regional demand accounting},
and \emph{commitment}.
A placement associates a request with a candidate serving replica and,
therefore, with the region containing that replica.
Regional demand accounting attributes an arrived request's prospective work
to its tentative replica and region so that the request influences subsequent
placement decisions.
Commitment is the boundary after which the routing layer no longer revises
that placement.
Using only causally available information, the router seeks to maximize the
fraction of requests that satisfy all applicable request-level SLOs.
RAVEL chooses successful handoff to the serving engine as its commitment
boundary: a request remains reassignable while it is Router-owned and becomes
final only after handoff.
This is a deliberate control-plane boundary rather than an inherent
irreversibility property of LLM execution; it preserves placement flexibility
without requiring migration of engine-managed KV state.



\subsection{Request-Dependent Regional Feasibility}
\label{sec:problem:feasibility}

Prior SLO-aware serving systems show that whether a request can meet its
latency objective depends on both request-specific requirements and evolving
serving conditions~\cite{distserve,jitserve,quartz,andes}.
RAVEL makes this dependence explicit at the geographic placement layer:
SLO feasibility is specific to a request, candidate replica, and decision
time rather than a static property of a region.

Let $c$ denote a serving region, $\mathcal{R}_c$ the replicas in that region,
and $r\in\mathcal{R}_c$ a candidate replica.
For request $i$, let $\Phi_{i,r}(t)$ indicate whether replica $r$ is predicted,
using only information available at time $t$, to satisfy every SLO constraint
applicable to $i$.
This prediction depends on the request's SLO semantics and remaining slack,
its estimated service demand, the WAN path to region $c$, and the currently
observed serving state at replica $r$.
Because both serving state and remaining slack evolve over time,
$\Phi_{i,r}(t)$ is an online feasibility estimate rather than a guarantee:
future arrivals and prediction error may still change the realized outcome.
Section~\ref{sec:design:feasibility} describes how RAVEL constructs these
replica-level predictions.

We define the feasible-region set of request $i$ as
\begin{equation}
    \mathcal{F}_i(t)
    =
    \left\{
        c
        \;\middle|\;
        \exists r\in\mathcal{R}_c:
        \Phi_{i,r}(t)=1
    \right\}.
    \label{eq:problem:feasible-set}
\end{equation}
The time index captures two distinct sources of change.
First, new arrivals and accounting updates alter the serving state visible to
the router, so $\mathcal{F}_i(t)$ may expand or shrink as contention changes.
Second, even if network and serving conditions remain unchanged, elapsed time
consumes TTFT or completion slack and can therefore remove regions from the
feasible set.

The size and composition of $\mathcal{F}_i(t)$ characterize how substitutable
regional capacity is for request $i$.
A request with several feasible regions has multiple serving alternatives,
whereas a request with a small feasible set depends on comparatively scarce
regional capacity.
For example, if
$\mathcal{F}_1=\{A,B\}$ and $\mathcal{F}_2=\{A\}$,
assigning request 1 to region $A$ may consume the only feasible option
available to request 2, even though request 1 could instead execute in $B$.
Thus, the opportunity cost of assigning capacity in a region depends on which
requests can substitute away from that region.
Regional capacity is therefore request-dependent, and optimizing
request--region assignments independently need not maximize aggregate
request-level SLO attainment.


\subsection{The Information--Slack Tradeoff}
\label{sec:problem:tradeoff}

The allocation problem is online: the router must make decisions without
observing future arrivals.
A placement that is reasonable under the currently visible state may become
undesirable after a more region-constrained request arrives.
This setting shares concerns with low-latency online and shared-state cluster
scheduling, where decisions must be made from incomplete and continually
changing information
~\cite{sparrow,flow_imprecise,omega,firmament}.
RAVEL does not assume an arrival oracle or attempt to reproduce an offline
optimal allocation; instead, it preserves bounded revisability using only
causally observed requests and serving state.

The relevant operating points differ along two independent axes:
when arrived demand becomes visible to capacity planning and when placement
becomes final.
Table~\ref{tab:problem:operating-points} separates these decisions.
A one-shot operating point assigns and accounts for a request immediately and
does not subsequently revise that placement.
At the other extreme, our diagnostic LateBind-$\tau$ policy keeps a request
centrally unassigned and unaccounted for a fixed interval $\tau$, then observes
the current state and makes a single final placement.
RAVEL occupies the remaining operating point: it makes the request's demand
visible immediately while retaining placement revisability until engine
handoff.

\begin{table}[t]
    \centering
    \small
    \caption{
    Demand-visibility and placement-finality operating points.
    One-shot couples both decisions at arrival, LateBind-$\tau$ delays both,
    and RAVEL exposes demand immediately while deferring only commitment.
    }
    \label{tab:problem:operating-points}

    \renewcommand{\arraystretch}{1.25}

    \begin{tabular*}{\columnwidth}{
        @{\extracolsep{\fill}}
        lcc
        @{}
    }
        \toprule
        \textbf{Policy} &
        \makecell{\textbf{Demand}\\\textbf{visible}} &
        \makecell{\textbf{Placement}\\\textbf{final}} \\
        \midrule
        One-shot        & At arrival   & At routing        \\
        LateBind-$\tau$ & After $\tau$ & After $\tau$      \\
        RAVEL           & At arrival   & At engine handoff \\
        \bottomrule
    \end{tabular*}
\end{table}

Deferring commitment can reveal information about later demand, but that
information is not free under request-level SLOs.
Delay scheduling similarly shows that waiting can expose better placement
opportunities~\cite{delayscheduling}, while deadline-driven scheduling
emphasizes the cost of consuming finite execution slack
~\cite{liu_layland_edf}.
Tail-sensitive distributed services further motivate treating latency slack
as a scarce resource rather than an unlimited observation window
~\cite{tail_at_scale}.
In geo-distributed LLM serving, these effects interact directly:
holding request $i$ for $\tau$ consumes $\tau$ of its available TTFT or
completion slack, leaving less time for WAN transfer, queueing, and inference.
A longer delay may therefore reveal more subsequent demand while
simultaneously reducing the set of regions in which the held request can still
succeed.
Equivalently, $\mathcal{F}_i(t)$ can shrink purely because the router waits,
even without an increase in serving load.

We refer to this tension as the \emph{information--slack tradeoff}.
Figure~\ref{fig:motivation} illustrates it empirically:
short fixed deferral can improve SLO attainment by exposing useful arrivals,
but the benefit saturates and eventually reverses as waiting consumes an
increasing fraction of request SLO budgets.
The design opportunity is therefore not simply to delay placement longer.
It is to make arrived demand visible immediately while deferring only
placement finality.

Exploiting this opportunity requires preserving useful revisability without
introducing unbounded waiting, inconsistent accounting, or unsafe ownership
transitions.
A practical design must
(i) estimate request-specific SLO feasibility,
(ii) account each tentative request exactly once and transfer its prospective
charge when placement changes,
(iii) bound both deliberate observation and multi-request replanning, and
(iv) enforce a failure-safe commitment boundary at engine handoff.
Section~\ref{sec:design} describes how RAVEL satisfies these requirements.
```




